% !TeX root = ../main.tex
\chapter{Założenia projektowe i analiza wymagań}
\label{chap:requirements}

% FIRST DRAFT. Requirements, actors, use cases and the class model come from
% the PAI course documents (SRODA4.*, October 2026); the diagrams are copied
% from them. Event dates come from collectors/config/events.yaml and record
% counts from the reporting views.

Ten rozdział opisuje, co system ma robić, dla kogo i na jakich danych.
Zaczynamy od studium przypadku, bo to ono wyznaczyło większość wymagań: liczbę i rodzaj źródeł, zakres czasu i skalę danych.
Dalej opisujemy źródła opinii, użytkowników systemu, wymagania funkcjonalne i niefunkcjonalne, a na końcu model przypadków użycia i model klas dziedziny.
Wymagania i modele opracowaliśmy jesienią 2026~r. w formie dokumentacji projektowej, a tutaj przedstawiamy je w skróconej postaci.

\section{Studium przypadku i oś wydarzeń}
\label{sec:req-case-study}

Studium przypadku jest gra \enquote{Clair Obscur: Expedition 33} francuskiego studia Sandfall Interactive.
\todo{Źródło: studio, wydawca, data premiery: find-sources.}
Wybraliśmy ją, bo jej historia nie jest płaska.
Grę zapowiedziano w czerwcu 2024~r., wydano 24 kwietnia 2025~r., kilka dni po premierze wybuchł spór o tekstury wygenerowane przez sztuczną inteligencję, potem przyszła fala uznania krytyków, a w grudniu 2025~r. seria nagród na gali The Game Awards i, dziewięć dni później, odebranie wyróżnień przez organizatorów Indie Game Awards.
Dla systemu, który ma wykrywać zmiany nastrojów, taka historia jest lepszym testem niż premiera bez niespodzianek.

Cykl życia gry opisaliśmy dwunastoma wydarzeniami, które przedstawia rysunek~\ref{fig:event-timeline}.
Daty wydarzeń zweryfikowaliśmy 12 września 2026~r. z publicznymi źródłami.
Dwa wpisy z pierwszej, roboczej listy okazały się błędne: pokaz, który uznaliśmy za prezentację Sony, był w rzeczywistości pokazem Xbox Developer Direct, a zapowiadane demo gry nigdy się nie ukazało.
Trzeci wpis, nominacja do nagrody The Game Awards 2024, nie miał miejsca i usunęliśmy go z listy.

\begin{figure}[hbp]
  \centering
  \begin{tikzpicture}[font=\footnotesize, x=1cm, y=1cm]
    % Steam reviews exist only from the launch onward.
    \fill[gray!30] (-0.12,-5.38) rectangle (0.12,-11.8);
    \draw[thick, -{Latex}] (0,0.3) -- (0,-12.0);
    \foreach \y/\t in {0/{VI 2024}, -3.5/{I 2025}, -9.5/{I 2026}} {
      \draw (-0.18,\y) -- (0.18,\y);
      \node[anchor=east] at (-0.25,\y) {\t};
    }
    \node[anchor=east, align=right] at (-0.25,-7.6) {okres recenzji\\Steam};
    \foreach \py/\ly/\txt in {
        -0.13/-0.13/{9.06.2024: zapowiedź na Xbox Games Showcase},
        -1.32/-1.12/{20.08.2024: pokaz dla prasy na Gamescom 2024},
        -1.43/-1.58/{27.08.2024: zwiastun rozgrywki \enquote{Flying Waters}},
        -3.87/-3.87/{23.01.2025: pokaz na Xbox Developer Direct 2025},
        -5.37/-4.62/{23.04.2025: koniec embarga na recenzje},
        -5.38/-5.06/{24.04.2025: premiera},
        -5.42/-5.50/{26.04.2025: spór o tekstury wygenerowane przez AI},
        -5.43/-5.94/{27.04.2025: milion sprzedanych egzemplarzy},
        -5.68/-6.38/{12.05.2025: fala uznania krytyków},
        -9.17/-8.98/{11.12.2025: The Game Awards 2025},
        -9.32/-9.44/{20.12.2025: Indie Game Awards odbierają wyróżnienia},
        -11.38/-11.38/{24.04.2026: aktualizacja rocznicowa}} {
      \fill (0,\py) circle (2pt);
      \draw[gray] (0.05,\py) -- (0.6,\ly);
      \node[anchor=west] at (0.65,\ly) {\txt};
    }
  \end{tikzpicture}
  \caption{Oś dwunastu wydarzeń cyklu życia gry; odległości na osi są proporcjonalne do czasu, szary pas oznacza okres, z którego pochodzą recenzje Steam}
  \label{fig:event-timeline}
\end{figure}

Wydarzenia zapisaliśmy w pliku konfiguracyjnym \code{collectors/config/events.yaml}.
Każde ma stały identyfikator, na przykład \code{launch} albo \code{tga_2025_sweep}, nazwę, typ i datę.
Identyfikator wydarzenia trafia do każdej opinii jako jej znacznik i jest podstawą wszystkich zestawień w podziale na wydarzenia.
Tabela~\ref{tab:event-coverage} pokazuje, ile opinii z każdego źródła ostatecznie przypisaliśmy do poszczególnych wydarzeń.

\begin{table}[hbp]
  \caption{Liczba opinii przypisanych do wydarzeń w podziale na źródła}
  \label{tab:event-coverage}
  \centering
  \small
  \begin{tabular}{|l|l|r|r|}
    \hline
    Data       & Wydarzenie                                  & YouTube       & Steam          \\
    \hline
    09.06.2024 & zapowiedź na Xbox Games Showcase            & \num{1711}    & --             \\
    20.08.2024 & pokaz dla prasy na Gamescom 2024            & --            & --             \\
    27.08.2024 & zwiastun \enquote{Flying Waters}            & \num{1257}    & --             \\
    23.01.2025 & pokaz na Xbox Developer Direct 2025         & \num{292}     & --             \\
    23.04.2025 & koniec embarga na recenzje                  & --            & --             \\
    24.04.2025 & premiera                                    & \num{909}     & \num{3915}     \\
    26.04.2025 & spór o tekstury generowane przez AI         & --            & \num{5742}     \\
    27.04.2025 & milion sprzedanych egzemplarzy              & \num{477}     & \num{12163}    \\
    12.05.2025 & fala uznania krytyków                       & \num{351}     & \num{13715}    \\
    11.12.2025 & The Game Awards 2025                        & \num{3328}    & \num{4749}     \\
    20.12.2025 & Indie Game Awards odbierają wyróżnienia     & \num{3502}    & \num{5075}     \\
    24.04.2026 & aktualizacja rocznicowa                     & \num{427}     & \num{895}      \\
    \hline
    --         & poza oknami wydarzeń                        & --            & \num{205385}   \\
    \hline
    \multicolumn{2}{|l|}{Razem}                              & \num{12254}   & \num{251639}   \\
    \hline
  \end{tabular}
\end{table}

Z tabeli widać asymetrię, która jest własnością źródeł, a nie błędem zbierania.
Recenzje w sklepie mogą pisać tylko osoby, które grę kupiły, więc przed premierą Steam milczy, a cała wiedza o nastrojach pochodzi z komentarzy pod zwiastunami.
Po premierze proporcje się odwracają.
Dwa wydarzenia nie mają żadnych danych: pokaz dla prasy na Gamescom nie miał publicznego materiału wideo z komentarzami, a koniec embarga na recenzje wypadł dzień przed premierą, więc recenzje z tego okresu przypisały się do premiery jako wydarzenia najbliższego w czasie.
\todo{Potwierdzić przyczynę braku danych dla Gamescom 2024.}

\section{Źródła danych}
\label{sec:req-sources}

Pierwotny projekt zakładał dwa źródła: komentarze z serwisu YouTube i wpisy z serwisu Reddit.
Dwa źródła nie były wyborem przypadkowym.
Głównym argumentem za odtwarzaniem strumienia w globalnej kolejności chronologicznej jest to, że opinie przychodzą z niezależnych serwisów, z różnymi zegarami; przy jednym źródle ten mechanizm działałby, ale niczego by nie dowodził.

\subsection{YouTube}
\label{sec:req-youtube}

Komentarze pobieramy przez YouTube Data API w wersji~3, metodą \code{commentThreads.list}, która zwraca komentarze najwyższego poziomu wraz z odpowiedziami~\cite{docyoutubecommentthreads}.
Każde wywołanie tej metody kosztuje jedną jednostkę dziennego limitu, który wynosi \num{10000} jednostek, a wyszukiwanie filmów metodą \code{search.list} kosztuje sto jednostek~\cite{docyoutubequota}.
Ta różnica przesądziła o sposobie doboru materiału.
Zamiast wyszukiwać filmy automatycznie, utrzymujemy ręcznie dobraną listę 23 filmów przypisanych do wydarzeń: oficjalnych zwiastunów, materiałów wydawcy i najczęściej komentowanych recenzji.
Każdy identyfikator filmu przed dodaniem do listy sprawdziliśmy przez API, łącznie z kanałem, datą publikacji i liczbą komentarzy.
Kolektor pobiera komentarze spod filmów przypisanych do wydarzenia i zostawia tylko te, które opublikowano w zadanym oknie czasowym.

\subsection{Reddit: odmowa dostępu}
\label{sec:req-reddit}

Kolektor serwisu Reddit napisaliśmy w czerwcu 2026~r., razem z kolektorem YouTube, zgodnie z dokumentacją uwierzytelniania OAuth tego serwisu.
Wniosek o dostęp do interfejsu programistycznego został jednak odrzucony i nie istniała żadna zgodna z regulaminem droga obejścia tej decyzji~\cite{redditdataapi}.
Pobieranie danych z pominięciem oficjalnego interfejsu odrzuciliśmy, bo narusza warunki korzystania z serwisu, a wynik pracy, którego nie da się powtórzyć zgodnie z prawem, nie jest wynikiem.
Rozważaliśmy też serwis X, ale jego bezpłatny dostęp do interfejsu nie pozwalał na pobieranie danych w potrzebnej skali.
Kod kolektora Reddit pozostał w repozytorium, kompletny i przetestowany, jako dowód wykonanej pracy i gotowe rozwiązanie na wypadek, gdyby dostęp kiedyś się pojawił.

\subsection{Steam}
\label{sec:req-steam}

Drugim źródłem zostały recenzje gry w sklepie Steam, co zapisaliśmy w decyzji architektonicznej z 20 września 2026~r.
Punkt dostępowy sklepu, opisany w dokumentacji Steamworks~\cite{docsteamreviews}, zwraca recenzje w formacie \ac{json} przez zwykłe połączenie HTTPS, bez klucza, bez rejestracji i bez wniosku o przydział limitu.
Nie ma tu więc procesu akceptacji, w którym ktoś mógłby nam odmówić.
Recenzje przychodzą stronami po sto, od najnowszych, a przy sprawdzeniu 5 września 2026~r. gra miała ich \num{277523}.

Steam okazał się lepszy od pierwotnego planu w trzech sprawach.
Schemat danych pasował bez zmian do naszego formatu rekordu.
Źródła się uzupełniały: YouTube obejmuje cały cykl życia, łącznie z okresem przed premierą, a Steam okres od premiery, ale ze znacznie większą liczbą opinii.
Najważniejsze było jednak pole \code{voted_up}, czyli deklaracja autora, czy poleca grę.
Daje ono słabą etykietę dla każdej recenzji, której ani YouTube, ani Reddit nie oferowały.

\subsection{Wspólny format rekordu}
\label{sec:req-record}

Wszystkie kolektory zapisują opinie w jednym formacie, przedstawionym w tabeli~\ref{tab:record-schema}.
Dzięki temu dalsze etapy przetwarzania nie wiedzą i nie muszą wiedzieć, z którego serwisu pochodzi rekord, a dodanie kolejnego źródła wymaga tylko nowego kolektora.

\begin{table}[hbp]
  \caption{Wspólny format rekordu opinii zapisywanego przez kolektory}
  \label{tab:record-schema}
  \centering
  \small
  \begin{tabular}{|l|p{0.62\textwidth}|}
    \hline
    Pole                    & Znaczenie                                                                                               \\
    \hline
    \texttt{id}             & identyfikator z przedrostkiem źródła, np. \texttt{youtube:Ugx...} lub \texttt{steam:<recommendationid>} \\
    \hline
    \texttt{source}         & źródło: \texttt{youtube} albo \texttt{steam}                                                           \\
    \hline
    \texttt{parent\_id}     & identyfikator komentarza nadrzędnego dla odpowiedzi; dla recenzji zawsze pusty                         \\
    \hline
    \texttt{created\_utc}   & chwila opublikowania opinii przez autora (UTC)                                                          \\
    \hline
    \texttt{collected\_at}  & chwila pobrania opinii przez kolektor (UTC)                                                             \\
    \hline
    \texttt{author\_hash}   & pseudonim autora: 16 znaków skrótu SHA-256 identyfikatora z tajną solą                                  \\
    \hline
    \texttt{text}           & treść opinii                                                                                            \\
    \hline
    \texttt{language}       & język wstępnie określony przez kolektor; potok nadpisuje go własnym wynikiem                            \\
    \hline
    \texttt{score}          & liczba polubień komentarza albo głosów \enquote{przydatna} dla recenzji                                 \\
    \hline
    \texttt{context\_id}    & identyfikator filmu YouTube albo numer aplikacji w sklepie Steam                                        \\
    \hline
    \texttt{event\_tag}     & identyfikator wydarzenia albo \texttt{general} dla recenzji spoza okien wydarzeń                        \\
    \hline
    \texttt{voted\_up}      & rekomendacja autora (tylko Steam); dla YouTube wartość pusta, nigdy fałsz                               \\
    \hline
  \end{tabular}
\end{table}

\section{Interesariusze i użytkownicy}
\label{sec:req-users}

Interesariuszami projektu są uczelnia, w imieniu której zakres prac zatwierdza promotor, przyszli odbiorcy wyników, zespół projektowy oraz, jako źródła wymagań nieosobowych, akty prawne i regulaminy serwisów, z których pochodzą dane.
Z rozmów z promotorem wynikły trzy cele biznesowe.
Pierwszym jest czytelny obraz zmian nastrojów wokół najważniejszych wydarzeń, bez ręcznego przeglądania dziesiątek tysięcy komentarzy.
Drugim jest możliwość wielokrotnego użycia systemu przy niskim koszcie: kolejne studium przypadku ma wymagać zmiany konfiguracji, a koszt chmury ma mieścić się w budżecie.
Trzecim jest wiarygodność wyników: jakość modelu ma być zmierzona, a ponowne przetworzenie tych samych danych ma dawać ten sam wynik.

System ma trzy role użytkowników.
Analityk jest odbiorcą wyników, na przykład pracownikiem działu marketingu wydawcy albo badaczem; korzysta wyłącznie z panelu i nie musi znać budowy systemu.
Operator systemu jest członkiem zespołu, który uruchamia pobieranie danych i odtworzenia, zatwierdza zmiany infrastruktury i reaguje na alarmy.
Specjalista uczenia maszynowego odpowiada za zbiór wzorcowy, trenowanie, ocenę i wdrażanie kolejnych wersji modelu.
Poza ludźmi system współpracuje z czterema systemami zewnętrznymi: YouTube Data API, serwisem recenzji Steam, narzędziem Data Studio i serwisem GitHub, w którym przechowujemy kod i który uruchamia automatyczne sprawdzanie zmian.

\section{Wymagania funkcjonalne}
\label{sec:req-functional}

Dwadzieścia wymagań funkcjonalnych podzieliliśmy na pięć grup, odpowiadających podsystemom.
Zestawia je tabela~\ref{tab:functional-requirements}; pełne opisy wraz z priorytetami i powiązaniami zawiera specyfikacja wymagań w dokumentacji projektu.

\begin{table}[hbp]
  \caption{Wymagania funkcjonalne w podziale na podsystemy}
  \label{tab:functional-requirements}
  \centering
  \small
  \begin{tabular}{|p{0.2\textwidth}|p{0.7\textwidth}|}
    \hline
    Podsystem                   & Wymagania                                                                                                                                                                                                                                                                      \\
    \hline
    Pozyskiwanie danych         & pobieranie komentarzy z YouTube dla wydarzenia i okna czasowego; pobieranie recenzji ze Steam z przypisaniem do najbliższego wydarzenia w promieniu 3 dni; pseudonimizacja autorów; zapis paczek w niezmiennym archiwum; konfiguracja studium przypadku w plikach              \\
    \hline
    Przetwarzanie strumieniowe  & ładowanie archiwum do hurtowni z usunięciem duplikatów; chronologiczne odtwarzanie strumienia z przyspieszeniem czasu; walidacja komunikatów i odkładanie błędnych do osobnej kolejki; deterministyczne wykrywanie języka i klasyfikacja wydźwięku; scalenie wyników z weryfikacją kompletności \\
    \hline
    Uczenie maszynowe           & trenowanie modelu z rejestracją przebiegów; ocena jakości na zbiorze wzorcowym i na ocenach autorów recenzji; rejestracja wersji i wdrożenie na czas przetwarzania                                                                                                        \\
    \hline
    Prezentacja wyników         & autoryzowane widoki raportowe bez danych osobowych; wykres zmian wydźwięku w czasie z datami wydarzeń; zestawienie dla wydarzeń                                                                                                                                         \\
    \hline
    Infrastruktura i utrzymanie & tworzenie zasobów wyłącznie z kodu; automatyczne sprawdzanie i wdrażanie zmian; odtworzenie uruchamiane jednym poleceniem; powiadomienia o awariach i przekroczeniu budżetu                                                                                              \\
    \hline
  \end{tabular}
\end{table}

Dwa wymagania zasługują na komentarz, bo zdecydowały o architekturze.
Wymaganie odtwarzania strumienia mówi, że rekordy z obu źródeł trafiają do kolejki w jednej, wspólnej kolejności czasu publikacji, z odstępami skróconymi przez współczynnik przyspieszenia, domyślnie tak, by jeden dzień rzeczywisty trwał jedną sekundę.
Wymaganie scalenia mówi, że wyniki przetwarzania trafiają do docelowej tabeli dopiero wtedy, gdy przetworzono wszystkie rekordy danego odtworzenia.
To drugie wymaganie dopisaliśmy po tym, jak jedno z wczesnych odtworzeń scaliło niepełny zbiór wyników, o czym piszemy w podrozdziale~\ref{sec:impl-replay-build}.

\section{Wymagania niefunkcjonalne}
\label{sec:req-nonfunctional}

Wymagania niefunkcjonalne zestawia tabela~\ref{tab:nonfunctional-requirements}.
Najważniejsze z nich dotyczą prywatności i powtarzalności, bo od nich zależy, czy wyniki w ogóle nadają się do pracy dyplomowej.

\begin{table}[hbp]
  \caption{Najważniejsze wymagania niefunkcjonalne}
  \label{tab:nonfunctional-requirements}
  \centering
  \small
  \begin{tabular}{|p{0.22\textwidth}|p{0.68\textwidth}|}
    \hline
    Wymaganie                    & Treść                                                                                                                                             \\
    \hline
    Ochrona prywatności autorów  & nazwy użytkowników nie trafiają do żadnego zapisu w chmurze; pseudonim autora i treść opinii nie opuszczają zbioru analitycznego; dane pozostają w regionie UE \\
    \hline
    Powtarzalność wyników        & dwa odtworzenia tych samych danych z tym samym klasyfikatorem dają identyczny zbiór wierszy, z pominięciem wersji modelu i czasu przetworzenia      \\
    \hline
    Kompletność danych           & każdy rekord trafia do tabeli zdarzeń albo do kolejki błędnych komunikatów; plików archiwum nie modyfikuje się ani nie usuwa                         \\
    \hline
    Bezpieczeństwo dostępu       & osobne konto usługi z minimalnymi uprawnieniami dla każdego komponentu; sekrety tylko w Secret Manager; maszyny bez publicznych adresów IP            \\
    \hline
    Kontrola jakości modelu      & nowa wersja zastępuje dotychczasową tylko przy wyższym F1 makro na części kontrolnej zbioru wzorcowego i nie gorszym F1 dla klasy negatywnej          \\
    \hline
    Wydajność                    & zbiór rzędu \num{300000} opinii przetwarzany przy najwyżej 2 maszynach roboczych, pełne odtworzenie w nie więcej niż 6 godzin                       \\
    \hline
    Limity serwisów              & pobieranie z YouTube w dziennym limicie \num{10000} jednostek; zapytania do Steam pojedynczo, z przerwą \qty{0.5}{\second}                          \\
    \hline
    Niezależność od premiery     & nowe studium przypadku bez zmian w kodzie, przez zmianę przedrostka nazw zasobów i plików konfiguracji                                             \\
    \hline
    Wymienność modelu            & model można wymienić bez zmian w potoku dzięki wspólnemu interfejsowi klasyfikatora                                                                 \\
    \hline
    Ograniczenie kosztów         & koszt miesięczny w budżecie; zasoby rozliczane w sposób ciągły działają tylko na czas przetwarzania                                                 \\
    \hline
  \end{tabular}
\end{table}

Ze specyfikacji pochodzą też cztery kryteria akceptacji.
Pierwszym jest pokaz pełnego przebiegu dla obu źródeł, od pobrania do panelu, w którym liczba identyfikatorów w tabeli zdarzeń zgadza się z liczbą w tabeli pośredniej pomniejszoną o błędne komunikaty.
Drugim jest test powtarzalności: dwa odtworzenia z deterministycznym klasyfikatorem mają dać identyczny odcisk tabeli, a przy prawdziwym modelu etykieta może zmienić się dla nie więcej niż \qty{0.1}{\percent} opinii.
Trzecim jest raport jakości modelu na części kontrolnej zbioru wzorcowego i na ocenach autorów recenzji Steam.
Czwartym jest zgodność infrastruktury z kodem: plan zmian Terraform nie wykazuje różnic względem działającego środowiska.
Do każdego z tych kryteriów wracamy w rozdziale~\ref{chap:evaluation}.

\section{Model przypadków użycia}
\label{sec:req-use-cases}

Wymagania funkcjonalne sformalizowaliśmy w postaci 33 przypadków użycia w ujednoliconym języku modelowania (\ac{uml}), po jednym diagramie dla każdego podsystemu.
Aktorami są trzy role użytkowników i cztery systemy zewnętrzne.
Dodaliśmy też dwóch aktorów uogólnionych: członka zespołu projektowego, który grupuje przypadki wspólne dla operatora i specjalisty uczenia maszynowego, oraz źródło opinii, którego specjalizacjami są YouTube i Steam, a w przyszłości każde kolejne źródło.
Każde z dwudziestu wymagań funkcjonalnych realizuje co najmniej jeden przypadek użycia.

Najwięcej przypadków ma podsystem przetwarzania strumieniowego, przedstawiony na rysunku~\ref{fig:uc-stream}.
Jego centralnym przypadkiem jest uruchomienie odtworzenia, które zawiera publikację chronologicznego strumienia, przetworzenie komunikatów w potoku i scalenie wyników z weryfikacją kompletności.
Odrzucenie błędnego komunikatu rozszerza przetwarzanie, bo zachodzi tylko dla rekordów, których nie da się odczytać.

\begin{figure}[hbp]
  \centering
  \includegraphics[width=\textwidth]{uc-stream-processing}
  \caption{Diagram przypadków użycia podsystemu przetwarzania strumieniowego}
  \label{fig:uc-stream}
\end{figure}

Podsystem uczenia maszynowego, przedstawiony na rysunku~\ref{fig:uc-ml}, obejmuje pełny cykl życia modelu: oznaczenie zbioru wzorcowego, trenowanie, ocenę, wdrożenie wersji i przywrócenie poprzedniej, gdy nowa okaże się gorsza.
Przywrócenie wymaga współpracy dwóch ról, bo specjalista wdraża poprzednią wersję modelu, a operator ponownie odtwarza dane, co nadpisuje etykiety w tabeli zdarzeń.

\begin{figure}[hbp]
  \centering
  \includegraphics[width=\textwidth]{uc-machine-learning}
  \caption{Diagram przypadków użycia podsystemu uczenia maszynowego}
  \label{fig:uc-ml}
\end{figure}

Podsystem analityczny, przedstawiony na rysunku~\ref{fig:uc-analytics}, jest jedynym, z którego korzysta analityk.
Przypadek uogólniony przeglądu wyników ma trzy specjalizacje, odpowiadające trzem stronom panelu: zmianom wydźwięku w czasie, zestawieniu wydarzeń i pokryciu źródeł i języków.

\begin{figure}[hbp]
  \centering
  \includegraphics[width=0.85\textwidth]{uc-analytics}
  \caption{Diagram przypadków użycia podsystemu analitycznego}
  \label{fig:uc-analytics}
\end{figure}

\todo[inline]{Diagramy przypadków użycia dla podsystemów pozyskiwania danych i infrastruktury oraz trzy diagramy sekwencji są w dokumentacji PAI (model przypadków użycia); zdecydować, czy dodać je tu, czy w dodatku.}

\section{Model klas dziedziny}
\label{sec:req-class-model}

Model klas opisuje dziedzinę problemu, a nie kod.
Klasy wynikają z wymagań na dane i odpowiadają pojęciom, którymi posługujemy się w całej pracy.
Rysunek~\ref{fig:class-opinions} przedstawia opinie i ich kontekst.

\begin{figure}[hbp]
  \centering
  \includegraphics[width=0.95\textwidth]{class-opinions}
  \caption{Diagram klas: opinie i ich kontekst}
  \label{fig:class-opinions}
\end{figure}

Centralną klasą jest abstrakcyjna opinia, której specjalizacjami są komentarz YouTube i recenzja Steam.
Komentarz może być odpowiedzią na inny komentarz i zawsze należy do filmu, a recenzja nie ma struktury wątków, ma za to ocenę autora.
Opinia jest przypisana do dokładnie jednego wydarzenia, a każde wydarzenie należy do studium przypadku.
Autor występuje w modelu wyłącznie jako pseudonim; model z założenia nie ma miejsca na nazwę użytkownika.
Paczka archiwum grupuje opinie pobrane w jednym przebiegu kolektora dla jednego źródła i wydarzenia i nigdy nie jest modyfikowana.

Rysunek~\ref{fig:class-processing} przedstawia klasy związane z przetwarzaniem.
Wynik klasyfikacji łączy opinię z etykietą wydźwięku, pewnością modelu i wersją modelu, która tę etykietę nadała.
Wyników dla jednej opinii może być wiele, bo każde odtworzenie tworzy nowy, ale w tabeli docelowej zostaje zawsze ten z najpóźniejszego przetworzenia.
Odtworzenie przechodzi kolejne stany, od uruchomienia, przez publikację i opróżnianie potoku, po scalenie albo przerwanie; jego cykl życia opisujemy dokładniej w następnym rozdziale.

\begin{figure}[hbp]
  \centering
  \includegraphics[width=0.8\textwidth]{class-processing}
  \caption{Diagram klas: przetwarzanie opinii}
  \label{fig:class-processing}
\end{figure}

Trzecia grupa klas, nieprzedstawiona na rysunkach, dotyczy oceny modelu: zbiór wzorcowy składa się z oznaczeń wzorcowych przypisanych do części uczącej, strojenia lub kontrolnej, a raport oceny łączy wersję modelu ze zbiorem wzorcowym i zawiera dokładność, F1 makro oraz wyniki dla klas i języków.
Tak zdefiniowana dziedzina wyznacza, jakie dane system musi przechowywać i jakie operacje na nich wykonywać.
Następny rozdział pokazuje, jak przełożyliśmy ją na architekturę w chmurze.
